\section{结论与展望}
本文研究了临床长文本场景中的可信输出层建模问题，提出了一种基于 \encoder{} 的临床长文本可信预测框架 \method{}。该方法通过冻结长上下文编码器、在特征空间中联合建模预测均值与输入相关方差，并在训练后期引入加权最小二乘精修机制，将深度语义表示与结构化统计求解结合起来，从而把研究重点从“获得更强表示”推进到“构建更可靠输出”。

结合本文的阶段性实验结果，完整 \method{} 在 Phenotype 多标签任务上达到 0.786 的 weighted F1、0.804 的 micro F1 和 0.889 的 macro AUROC，同时在 MSE 和 MAE 上优于线性头、MLP 头、Ridge 头以及阶段 A 模型。这一结果趋势表明，当长文本编码器已经提供较强表示能力时，围绕输出层继续开展可信建模仍然具有明确研究价值；而阶段 B 相比阶段 A 的进一步提升，则说明基于样本权重的结构化精修并非附加装饰，而是方法有效性的组成部分。

本文工作的意义主要体现在两个方面。其一，本文说明了在不重新训练更大基础模型的条件下，仍可围绕输出层形成具有方法深度的研究路径。其二，本文提出的框架为后续扩展提供了清晰接口，可继续引入 Bayesian Last Layer、conformal prediction、局部参数高效微调以及更细致的临床误差分析，从而把当前工作进一步扩展为更完整的临床可信预测系统。

后续工作主要包括四个方向。第一，在统一的数据整理与训练脚本基础上完成完整复现实验，并通过多随机种子检验结论稳定性。第二，围绕 pooling 策略、损失函数和 WLS 正则参数开展更系统的消融分析。第三，在当前输出层基础上接入 Bayesian Last Layer 与 conformal prediction，使模型从“可输出方差”进一步走向“可校准、可拒识、可审计”。第四，探索部分解冻编码器或 PEFT 策略，研究在保持成本可控的前提下进一步提升表示质量的可能性。
